\documentclass[10pt]{article}
\usepackage[margin=20mm]{geometry}
\usepackage{amsmath,amssymb,fontspec,unicode-math,xurl,hyperref}
\IfFontExistsTF{DejaVu Serif}{\setmainfont{DejaVu Serif}}{\setmainfont{Latin Modern Roman}}
\setmathfont{latinmodern-math.otf}
\setlength{\parindent}{0pt}\setlength{\parskip}{6pt}\sloppy
\title{Initial Scientific Review of
Automatic Differentiation of Optimization Algorithms with Time-Varying Updates}
\author{Assigned review 0249 / I25-43618}
\date{Initial review; scientific adjudication pending}
\begin{document}
\maketitle
Abstract: The general rate bound for nonconvergent update mappings is false: individual Jacobian spectral radii do not control products of varying nonnormal matrices. A complete affine counterexample satisfies every stated assumption. The accelerated proximal-gradient theorem also fails under its printed hypotheses because the curvature conditions used by its proof are omitted; a smooth convex example has stationary primal iterates but quadratically divergent derivatives. Both are Major findings concerning central theoretical contributions. Neither establishes failure of the better-conditioned special cases or falsity of the reported experiments. The published experimental recipe has a separate parameter-domain inconsistency, and its actual run settings remain Needs Evidence.\par
\section*{Paper, version, contribution, and scope}
The reviewed original is Sheheryar Mehmood and Peter Ochs, ICML 2025, PMLR 267 [paper]. Its 22-page publisher PDF includes Appendices A–D; printed and PDF page numbers coincide. The original, derived text and saved publisher page match all three manifest-bound hashes. No verified author implementation is staged. This means source-level and historical-run checks are unavailable, not that no implementation exists. No author code, data download or benchmark rerun was performed.\par
The paper studies differentiation of \(x^{(k+1)}(u)=A_k(x^{(k)}(u),u)\), including sequences without a limiting update map. It seeks convergence and rates for the corresponding derivative recursion, then specializes to proximal-gradient descent (PGD) and accelerated PGD (APG) on partly smooth objectives. The review read the full original, traced the main rate dependencies through Appendices A–C, independently calculated the witnesses below, and checked the experimental recipe in Appendix D. Source pages 5, 9, 14, 21 and 22 were visually inspected. A bounded primary-source lookup found arXiv v2 dated 24 October 2024 and an author page with the same update date [preprint, authorpage]; no later correction was verified. This was a metadata/page check, not a complete preprint comparison.\par
\section*{Problems and independent checks}
Problem 1 (Major: individual spectral radii give a false general rate).  Theorem 2.7, Equation (7), p. 5, uses \(\rho=\limsup_k\rho(D_xA_k(x^*,u^*))\) in its error-rate bound. The invoked Theorem A.6, pp. 14–15, instead uses the limsup of an operator norm in one common norm. Appendix B.3, p. 18, does not justify this replacement.\par
Take \(\mathcal X=\mathcal U=\mathbb R^3\), Euclidean norms, \(x^*=u^*=0\), \(x^{(0)}=e_3\), and \(C=e_2e_1^\top\). Put \[B_{2j}=\begin{pmatrix}0&1/2&0\\0&0&0\\0&0&1/4\end{pmatrix},
 \qquad
 B_{2j+1}=\begin{pmatrix}0&0&0\\1/2&0&0\\0&0&1/4\end{pmatrix},
 \qquad A_k(x,u)=B_k(x-Cu)+Cu .\] Every map is affine and smooth on the common neighborhood \(\mathbb R^6\), and has the same unique fixed-point map \(\psi(u)=Cu\). Thus \(X^*=D\psi(0)=C\) and \[D_xA_k=B_k,\qquad D_uA_k=(I-B_k)C,\qquad
 B_kC+(I-B_k)C=C.\] The common Euclidean operator norm is \(\|B_k\|_2=1/2<1\), so Assumption 2.1(iii) holds. The primal iterates are exactly \(x^{(k)}(0)=4^{-k}e_3\), hence \(q_x=1/4\). All derivatives are constant in \((x,u)\), so the difference in Assumption 2.3 is identically zero. This checks smoothness, both fixed-point identities, contraction, neighborhood membership and the required primal/derivative-variation rate.\par
Nevertheless, every individual \(B_k\) has spectral radius \(1/4\), whereas the derivative error \(E_k=D x^{(k)}(0)-C\) satisfies \[E_0=-C,\qquad E_{k+1}=B_kE_k,\qquad
 \|E_k\|_2=2^{-k}.\] The alternating off-diagonal entries successively move the nonzero first column between \(e_1\) and \(e_2\), multiplying it by \(1/2\). Choosing \(\delta=1/8\), the printed theorem would require, for some fixed finite constants and all sufficiently large \(k\), \[2^{-k}\le [\xi_1+\xi_2(k-K)](3/8)^{k-K}.\] This is impossible: the ratio grows as \((4/3)^k\), faster than any linear factor. Reviewer-owned exact rational recurrence checks confirm these identities through 80 iterations; the exponential argument establishes the failure for arbitrarily large \(k\).\par
The core consequence is a false advertised rate in the new setting without convergent update maps. Qualitative convergence in Theorem 2.5 is consistent with this example: the derivatives do converge, but at rate \(1/2\), not the stated spectral-radius rate. A bound based on a common norm or on suitable matrix-product growth can avoid the counterexample. This is not a verified repair of every constant in the paper.\par
The counterexample should not be transferred automatically to structured PGD or convergent-Jacobian APG. For a quadratic proximal-gradient family with Hessians \(F,G\succeq0\) and \(H=F+G\succ0\), \[B_\alpha=(I+\alpha G)^{-1}(I-\alpha F)
          =I-\alpha(I+\alpha G)^{-1}H\] is self-adjoint in the common \(H\) inner product, since \(HB_\alpha\) is symmetric. That additional structure can make spectral radii control a common operator norm. A convergent matrix sequence also admits an adapted norm near its limiting matrix, as in Corollary A.7. These are substantive boundaries on Problem 1, not evidence that the general theorem is correct.\par
Problem 2 (Major: the printed APG theorem omits needed curvature hypotheses).  Theorem 3.12, p. 9, explicitly assumes 3.1 and 3.3, but not restricted positive definiteness in 3.2. Appendix C.4, p. 21, invokes Theorem 3.11, which requires 3.2(ii) for contraction and additionally 3.2(i) for the differentiable solution map.\par
Let \(\mathcal X=\mathcal U=M=\Omega=\mathbb R\), \(g=0\), and \[f(x,u)=\tfrac12\{x^2-\log(1+x^2)\}-ux,\qquad
 \nabla_x f=\frac{x^3}{1+x^2}-u,\qquad
 \nabla_x^2f=\frac{x^2(x^2+3)}{(1+x^2)^2}.\] This is a smooth convex objective with globally Lipschitz gradient of constant \(L=9/8\), as \[0\le \nabla_x^2f\le\frac98,\qquad
 \frac98-\nabla_x^2f=\frac{(x^2-3)^2}{8(1+x^2)^2}.\] The zero regularizer is partly smooth on \(\mathbb R\) with identity proximal map. At \(u^*=x^*=0\), non-degeneracy holds because \(\partial_x(f+g)=\{0\}\). Choose \[x^{(-1)}=x^{(0)}=0,\qquad \alpha_k=\tfrac12,\qquad
 \beta_k=\frac{k}{k+6}.\] The steps lie in a fixed compact subinterval of \((0,2/L)\); \(\beta_k\in[0,1]\) converges to one. Every primal iterate is exactly zero. Also \(D_xA_{\alpha_*}(0,0)=1>-1/3=-1/(1+2\beta_*)\), satisfying the theorem’s displayed spectral side condition. All additional local smoothness requirements hold.\par
Differentiating the actual recursion at \(u=0\) yields \[D_{k+1}=(1+\beta_k)D_k-\beta_kD_{k-1}+\tfrac12,\qquad
 D_{-1}=D_0=0,\qquad
 D_k=\frac{k(k+13)}{28}.\] The closed form follows by substitution, and exact rational checks through 100 iterations agree. The derivatives diverge quadratically even though the primal sequence is stationary. The unique minimizer solves \(\psi(u)^3/(1+\psi(u)^2)=u\). For \(u_n=1/[n(n^2+1)]\), it is \(\psi(u_n)=1/n\), giving \(\psi(u_n)/u_n=n^2+1\to\infty\). Thus a differentiable solution map cannot be inferred under the printed assumptions.\par
The nearby discussion describes identification and linear convergence under 3.1, 3.2(ii) and 3.3; it does not explicitly import 3.2 as a blanket assumption into all later statements. The preceding theorems list their curvature assumptions separately. Accordingly this finding concerns Theorem 3.12 as printed and the unsupported invocation in its proof. Restoring the omitted curvature assumptions excludes the witness. APG satisfying those stronger assumptions is not refuted, and the later rate clause with an empty interval when \(\rho=1\) does not rescue the theorem’s earlier unconditional derivative-convergence claim.\par
Problem 3 (Minor specification inconsistency; historical compliance Needs Evidence).  Appendix D, p. 22, specifies \(\lambda\sim N(0,10^{-3})\) for the logistic experiments, while requiring \(\lambda>0\) and calling the objective \(\lambda\)-strongly convex. An untruncated centered normal puts probability \(1/2\) on negative values. For negative \(\lambda\), its quadratic regularizer dominates the at-most-linear logistic-loss growth along any nonzero ray, making the objective unbounded below.\par
This is a demonstrated inconsistency in the published recipe. It is not evidence that actual runs used negative values: truncation, absolute values, rejection sampling or a transcription error could explain it. The actual draws and implementation are unavailable. Section 4 states that the assumptions were satisfied, which is relevant counterevidence about the intended experiment.\par
There is also a theorem/run correspondence gap. Under independent infinite sampling, the stated PGD\(_{k_1}\) law \(U(0,2/(3L))\) has infimum zero almost surely, and PGD\(_{k_3}\) with \(U(4/(3L),2/L)\) has supremum \(2/L\). These fail Algorithm 1’s uniform bounds \(0<\underline\alpha\le\alpha_k\le\bar\alpha<2/L\). Finite sampled trajectories can have valid interior minima/maxima and exhibit the plotted decay; their existence does not verify the stated asymptotic hypotheses. The sampling scheme, clipping and continuation convention are Needs Evidence.\par
\section*{Overall impact, positive checks, and state}
The overall assessment is Major, based independently on Problems 1 and 2. A wrong core rate and an APG convergence claim admitting divergent derivatives are material theoretical failures even without demonstrated benchmark-score loss. The fixed-point derivative recursion and the affine common-norm convergence mechanism pass the bounded checks above. The common-metric quadratic PGD identity and limiting-matrix norm argument provide concrete counterevidence against an indiscriminate rejection of every specialized result. They are not an exhaustive certification of Appendices A–C.\par
The empirical comparison uses five methods, 50 lasso instances, 50 perturbed logistic instances and one CIFAR10 problem. It warms starts within \(10^{-2}\) of a computed solution and compares a shared directional derivative per problem, rather than the full Jacobian norm. These are disclosed scope limits. Figure 1 shows the reported local behavior, but seeds, exact solver tolerances, raw error sequences, derivative directions and reference-solution errors were not independently verified. No issue is inferred solely from those unavailable materials.\par
Item | Assessment | Repair check | Adoption | Remaining work\par
P1 | Major | Common-norm boundary only | Unknown | Replace general spectral-radius rate and verify constants\par
P2 | Major | Added curvature excludes witness | Unknown | State hypotheses and recheck the complete APG proof\par
P3 | Minor | Domain contradiction verified | Unknown | Supply actual regularization draws and sampling rules\par
Scientific adjudication is pending; human reviews are empty. No acceptance decision or author adoption is asserted.\par
\section*{Supporting checks}
.\par
The source manifest and original bytes are unchanged. There is no verified code archive or historical execution identity to bind. Run reviewer-owned with Python; it uses exact arithmetic, imports no author module, and saves . The analytic proofs supply the infinite-iteration conclusions beyond those finite checks.\par
Only the five original pages specified above were visually inspected.\par
9 S. Mehmood and P. Ochs. Automatic Differentiation of Optimization Algorithms with Time-Varying Updates. PMLR 267 (2025). https://proceedings.mlr.press/v267/mehmood25a.html. S. Mehmood and P. Ochs. Preprint metadata, arXiv:2410.15923v2, 24 October 2024; accessed 25 October 2026. https://arxiv.org/abs/2410.15923. Authors’ publication page; accessed 25 October 2026. https://www.mop.uni-saarland.de/pub/MO25/index.shtml.\par
\end{document}
